\section{Limitations}

First, \core{} is a distilled benchmark rather than a replacement for the full 664-task reservoir. It is intended for rapid comparison and development, while full-reservoir evaluation remains useful for final audits.

Second, probe-based selection depends on the selected probe algorithms. \arena{} mitigates this with panel fidelity and complementarity, but future algorithm paradigms may require updating the probe set.

Third, the ground-truth focus favors equation-discovery evaluation. Real-world black-box regression tasks without known formulas are not the primary focus of \arena{}.

Fourth, LLM-assisted methods may benefit from external priors. \arena{} partially controls this by separating no-context probe evaluation from context-enabled full evaluation, but contamination and memorization remain important concerns.

Fifth, ROBU requires additional noisy-training runs. The robustness axis is meaningful only when all submitted algorithms run the same noisy-training protocol.

Finally, symbolic equivalence is imperfect. CAS simplification and numeric equivalence checks can fail on pathological expressions, so thresholds, probe domains, and failure cases must be documented.
